\cotefile{10}
\scanpage{42}
\mspage{\handwritten{10}}
\uline{pe de présentation finie, $H_j$ un $S_j$-préschéma de présentation finie à groupe d'opérate\deleted{ru}\inserted{ur}s $G_j$. Pour que $H=H_j\handwritten{\times}_{S_j}S$ soit un $S$-préschéma en espaces homogènes sous $\deleted{\unsure{H|G}}$}\margin{\uline{$G=G_j\handwritten{\times}_{S_j}S$}}\uline{ (resp. un fibré principal homogène sous $\deleted{G}\inserted{\deleted{G=G\unsure{?}S}}$\unsure{,|),} pour la }\uuline{topologie fidèlement plate localement de présentation finie}\uline{, il faut et il suffit qu'il existe un indice $i\handwritten{\in}I$, $i\handwritten{\geqslant}j$, tel que $H_i=H_j\handwritten{\times}_{S_j}S_i$ soit un $S_i$-préschéma en espaces homogènes sous $\deleted{G_{\unsure{i|j}}\unsure{?}}$}\margin{\uline{$G_i=G_j\handwritten{\times}_{S_j}S_i$}}\uline{ (resp. un fibré principal homogène sous $G_i$).}

Compte tenu de (10.14) et de (EGA IV 8.8.2, 8 8.3 et 8.10.5) , \deleted{\unsure{l'expr|l'exp}} l'énoncé résulte de l\deleted{j}\inserted{a} propriété concernant les morphismes \deleted{concernant}\inserted{couvrants} pour la topologie fppf rappelée en (10.0).

\figure{\begin{tikzpicture}[x=.5cm,y=.5cm]
\draw[line width=.3pt] (0,11.6) .. controls (3.4,11.6) and (6.5,11.1) .. (5.1,10.2)
 .. controls (3.7,9.3) and (-1.8,7.2) .. (-.1,7.1)
 .. controls (1.9,7.0) and (4.1,7.6) .. (3.7,5.8)
 .. controls (3.5,4.8) and (2.1,3.8) .. (2.7,2.9)
 .. controls (3.4,1.9) and (5.0,1.1) .. (4.5,0);
\draw[line width=.55pt] (2.45,11.7) .. controls (2.5,7.2) and (2.9,4.0) .. (3.75,.1);
\end{tikzpicture}}
% archive sheet number: 41
